% -----------------------------------------------------------------------------------
\section{Axioms for counterfactual operators}\label{sec:operators}

Section~\ref{sec:inference} leaves one question: is there a coherent way to assign outcomes to the
non-actual outputs of a computation? This is the problem of \emph{logical counterfactuals}. This
section states candidate desiderata precisely and determines every consistent and inconsistent
combination of them. There is a sharp impossibility among three individually satisfiable axioms,
and it identifies exactly which notion of program equivalence is responsible. Under a different
notion, all the desiderata hold together.

\subsection{Preliminary repairs}

Three features of a naive formulation produce failures before the substantive question is reached.
They are recorded here and excluded below by stipulation.

\emph{Transported contexts need not terminate.} Let a counterfactual operator be required to
\emph{transport} dependence: a component $C[T]$ that computes a context $C$ applied to the agent is
evaluated, under the alternative, as $C$ applied to the altered agent. It is natural to promise only
that the host environment terminates when its components are replaced by tables. That promise is too
weak. Let $T\equiv0$, let $C[P](x)$ return $0$ if $P(x)=0$ and diverge if $P(x)=1$, treat $g=C[T]$ as one
opaque component, and let the host return $g(d)$. The host terminates under every whole-component table
replacement, and $g$ is the total constant $0$, yet the transported execution $C[T^{d\leftarrow1}](d)$
diverges. On this domain totality and execution-faithful transport are already inconsistent. The
domain is therefore required to be \emph{intervention-total}: every supplied computational context
terminates under every permitted substitution, as can be ensured by finite tables, a terminating
language, certificates, or an explicit promise. Atomic boundaries matter, since replacing an opaque
outer component by a table can hide divergence inside it.

\emph{Source-changing replacement changes inputs.}

\begin{proposition}\label{prop:rawsource}
If the description supplied to the agent is the literal source and replacing a component changes it,
baseline correctness is inconsistent with invariance under replacement of a component by an
extensionally equal one.
\end{proposition}

\begin{proof}
Let $E$ ignore the agent's action and return $g(d(E))$ with $g(x)=|x|\bmod2$. Replacing $g$ by an
extensionally identical implementation padded by one character changes the parity of the source,
and hence the actual payoff, while a constant agent takes the same action in both versions.
\end{proof}

The components are equal but are called on different inputs. Below, comparisons of implementations
preserve the distinguished inputs, and counterfactual transformations retain the original
observation.

\emph{Single-input transport changes one input.} An intervention at $d$ replaces $T$ by
\[
  T^{d\leftarrow a}(x)=\begin{cases}a,&x=d,\\T(x),&x\ne d,\end{cases}
\]
and strong transport evaluates $C[T^{d\leftarrow a}]$. A context $C[T]=T\circ h$ does not in general
become the constant $a$: if $h(d)\neq d$, then $C[T^{d\leftarrow a}](d)=T(h(d))$. Replacing the whole agent
by a constant is a different, policy-level intervention.

\subsection{The finite context domain}

The main theorems use a finite domain in which termination and program equivalence raise no
difficulty, so that what remains is the semantic question alone.

\begin{definition}[Context environment]\label{def:context}
A \emph{context environment} is $E=(u;c_1,\dots,c_k)$, where the action set is $A=\{0,1\}$, each
component $c_i\colon A\to B_i$ is a table with values in a finite set $B_i$, and the host
$u\colon A\times\prod_iB_i\to\mathbb Q$ is a finite payoff table. The argument of $c_i$ is the
\emph{live parameter}: the agent's output at the distinguished input, as seen by component $i$. $E$ is
\emph{independent} if every $c_i$ is constant. Write
\[
  V_E(a)=u\bigl(a,c_1(a),\dots,c_k(a)\bigr),\qquad
  W_{E,b}(a)=u\bigl(a,c_1(b),\dots,c_k(b)\bigr).
\]
\end{definition}

$V_E$ is coherent substitution through every context; $W_{E,b}$ changes the action at the port and
freezes every component at its value under the baseline $b$. The definition of independence is not
circular: it is invariance under every substitution of the declared live parameter, not the absence
of an undefined relation of \enquote{depending on $T$}. The domain is the causal-matrix domain
with its factorisation made explicit.

\begin{proposition}\label{prop:contextmatrix}
Let $M_E[b,a]=W_{E,b}(a)$. Then $M_E$ is the causal matrix of the one-shot problem in which an
infallible predictor forecasts the agent's output $b$ and the components respond to the forecast, and
$M_E[a,a]=V_E(a)$. Every real $2\times2$ matrix with rational entries is $M_E$ for some context
environment, and $E$ is independent exactly when all rows of $M_E$ coincide.
\end{proposition}

\begin{proof}
With structural states $b\in A$, $\mu_a=\delta_a$, and $U(b,a)=u(a,c(b))$, the causal matrix is
$\sum_s\mu_b(s)U(s,a)=W_{E,b}(a)$. Given $M$, take one identity component $c(z)=z$ and $u(a,x)=M[x,a]$. If
every $c_i$ is constant, $W_{E,b}$ does not depend on $b$; if some row differs from another, some
component's value differs between the two baselines.
\end{proof}

Strictly, the converse in the last clause holds for the matrix, not for the environment: distinct
component values may be invisible to the host. Nothing below uses the converse.

\begin{definition}[Counterfactual kernel]\label{def:kernel}
A \emph{kernel} assigns to every context environment $E$, baseline $b\in A$, and action $a\in A$ a value
$\kappa(E,b,a)\in\mathbb Q\cup\{\bot\}$, where $\bot$ means undefined. The baseline is the agent's actual
output, $b=T(d(E))$. A deterministic agent $\tau$ on context environments is
\emph{$\kappa$-ratifiable} at $E$ if $\kappa(E,\tau(E),\cdot)$ is defined at every action and is
maximised at $\tau(E)$.
\end{definition}

That the kernel depends on $T$ only through the local baseline is an explicit stipulation;
the program extension in Section~\ref{sec:programmap} keeps the same interface. Agent extensionality
is then automatic, because agents computing the same function have the same baseline. Two
notions of extensionality for the environment must be distinguished.

\begin{definition}[Grounded and parametric equivalence]
$E\sim_bE'$ (\emph{grounded equivalence at $b$}) if $E$ and $E'$ have the same host and interfaces and
$c_i(b)=c'_i(b)$ for every $i$. $E\approx E'$ (\emph{parametric equivalence}) if they have the same
host and interfaces and $c_i=c'_i$ as functions on $A$.
\end{definition}

Grounded equivalence has a direct reading for programs. Once the actual agent is inserted, a
component is a closed computation returning $c_i(b)$, and two closed computations can be
extensionally equal although their open contexts would respond differently to another action. The
function $x\mapsto T(x)$ and the constant function with value $T(d)$ agree at $d$, and after
insertion of $T$ both are closed programs with output $T(d)$. Parametric equivalence identifies
contexts only when they agree under every substitution, and keeps the information that grounded
equivalence discards.

\begin{definition}[Desiderata]\label{def:desiderata}
A kernel may satisfy the following conditions, for all context environments $E$, $E'$ and all
$a,b\in A$.
\begin{description}[font=\normalfont\sffamily,leftmargin=3.2em,labelwidth=2.4em,labelsep=.8em,itemsep=.35em]
  \item[CF1] \emph{Baseline.} $\kappa(E,b,b)=V_E(b)$.
  \item[CF2] \emph{Totality.} $\kappa(E,b,a)\ne\bot$.
  \item[CF3] \emph{Independent calibration.} If $E$ is independent, $\kappa(E,b,a)=V_E(a)$.
  \item[CF4] \emph{Transport.} If $\kappa(E,b,a)$ is defined, it equals $V_E(a)$.
  \item[CF5] \emph{Grounded extensionality.} If $E\sim_bE'$, then $\kappa(E,b,a)=\kappa(E',b,a)$, including
        equality of definedness.
  \item[CF6] \emph{Conservativity.} Every defined value has the form $u(a,z_1,\dots,z_k)$ with $z_i=c_i(b)$
        for every constant context $c_i$.
  \item[CF7] \emph{Composition.} A transformed environment is the record $(E,b;a)$, whose evaluator
        returns $\kappa(E,b,a)$, and a further intervention overwrites the target:
        $(E,b;a)[a\leftarrow z]=(E,b;z)$.
  \item[CF8] \emph{Computability.} $\kappa$, with $\bot$ returned explicitly, is computable.
  \item[CF9] \emph{Realisability.} Every $\kappa(E,b,a)$ is defined and is the payoff that some agent taking
        $a$ obtains under coherent substitution of the live parameter.
  \item[CF10] \emph{Ratifiability.} Some computable $\tau$ is $\kappa$-ratifiable at every $E$.
\end{description}
\end{definition}

Three remarks fix the reading. \ax{CF4} transports the whole supplied context and preserves the
host, and it is conditional on definedness; a version requiring the transported values to exist
would incorporate \ax{CF2}. Every agent taking $a$ obtains exactly $V_E(a)$ under slot substitution,
and constant agents supply the alternatives, so
\[
  \ax{CF9}\iff\ax{CF2}+\ax{CF4}.
\]
And \ax{CF7} is deliberately weak: every kernel has such a lift, since the record exists even when its
evaluator returns $\bot$. A history-free requirement, or one demanding fresh rematching after each
intervention, would be a different axiom (Section~\ref{sec:candidates}). \ax{CF6} constrains only the
constant contexts; how the others respond is left open. On the finite domain \ax{CF8} holds for every
kernel defined by a finite recipe from the tables, and \ax{CF7} for every kernel; both are recorded
because they become substantive on program domains.

\subsection{Four kernels}

\begin{definition}[Witness kernels]
\[
  \kappa_{\mathrm{slot}}(E,b,a)=V_E(a),\qquad
  \kappa_{\mathrm{port}}(E,b,a)=W_{E,b}(a),
\]
\[
  \kappa_{\mathrm{base}}(E,b,a)=\begin{cases}V_E(b),&a=b,\\\bot,&a\ne b,\end{cases}\qquad
  \kappa_{\mathrm{flat}}(E,b,a)=V_E(b).
\]
\end{definition}

By Proposition~\ref{prop:contextmatrix}, $\kappa_{\mathrm{slot}}$ reads the diagonal of the causal
matrix and $\kappa_{\mathrm{port}}$ reads row $b$: they are ex-ante and causal evaluation, now
appearing as counterfactual semantics. $\kappa_{\mathrm{base}}$ values only the actual output, as
partial semantics do; $\kappa_{\mathrm{flat}}$ assigns the actual payoff to every alternative, so
every agent is ratifiable under it. All four satisfy \ax{CF1}, \ax{CF7}, and \ax{CF8}. Their other
properties are as follows.

\begin{lemma}\label{lem:witnesses}
$\kappa_{\mathrm{slot}}$ satisfies every desideratum except \ax{CF5}, and the agent choosing the first
maximiser of $V_E$ ratifies it. $\kappa_{\mathrm{port}}$ satisfies \ax{CF2}, \ax{CF3}, \ax{CF5}, and
\ax{CF6}. $\kappa_{\mathrm{base}}$ satisfies \ax{CF4}, \ax{CF5}, and \ax{CF6}. $\kappa_{\mathrm{flat}}$
satisfies \ax{CF2}, \ax{CF5}, and \ax{CF10}.
\end{lemma}

\begin{proof}
$\kappa_{\mathrm{slot}}$ is total and equals $V_E$, which gives \ax{CF2}--\ax{CF4} and \ax{CF9}; its
values $u(a,c(a))$ have $c_i(a)=c_i(b)$ for constant contexts, giving \ax{CF6}; the first maximiser of
$V_E$ is computable from the tables. $\kappa_{\mathrm{port}}$, $\kappa_{\mathrm{base}}$, and
$\kappa_{\mathrm{flat}}$ depend on $E$ only through $u$ and the baseline outputs $c_i(b)$, which grounded
equivalence preserves, so all three satisfy \ax{CF5}. $\kappa_{\mathrm{port}}$ is total; on an
independent environment $W_{E,b}=V_E$; and its values have the form required by \ax{CF6}.
$\kappa_{\mathrm{base}}$ is defined only at $a=b$, where its value $V_E(b)=u(b,c(b))$ satisfies \ax{CF4}
and \ax{CF6}. $\kappa_{\mathrm{flat}}$ is total and constant in $a$, so every agent ratifies it.
\end{proof}

Nothing more holds: $\kappa_{\mathrm{port}}$ violates \ax{CF4}, \ax{CF9}, and \ax{CF10} on the
environments of Theorems~\ref{thm:transport} and~\ref{thm:calibration} below; $\kappa_{\mathrm{base}}$ is
not total and so violates \ax{CF2}, \ax{CF3}, \ax{CF9}, and \ax{CF10}; and $\kappa_{\mathrm{flat}}$
violates \ax{CF3}, \ax{CF4}, \ax{CF6}, and \ax{CF9} on any independent environment whose two
actions have different payoffs.

\subsection{The projection lemma and two theorems}

For a context environment $E$ and a baseline $b$, let $E^b$ replace every component by the constant
with its baseline value, $c^b_i(a)=c_i(b)$. Then $E\sim_bE^b$, $E^b$ is independent, and
$V_{E^b}(a)=W_{E,b}(a)$.

\begin{lemma}[Projection]\label{lem:projection}
\ax{CF3} and \ax{CF5} together imply $\kappa=\kappa_{\mathrm{port}}$. \ax{CF5} and \ax{CF6} together
imply $\kappa(E,b,a)=W_{E,b}(a)$ whenever $\kappa(E,b,a)$ is defined.
\end{lemma}

\begin{proof}
\ax{CF3} defines the values at $E^b$ as $V_{E^b}(a)=W_{E,b}(a)$, and \ax{CF5} transfers them, with
their definedness, to $E$. For the second claim, transfer a defined value to $E^b$ by \ax{CF5}. Every
component of $E^b$ is constant, so \ax{CF6} fixes every component value at $c_i(b)$, and the value is
$u(a,c(b))=W_{E,b}(a)$.
\end{proof}

This is the main structural fact: \emph{calibration together with the erasure of dependency
information forces causal evaluation}. The two semantic requirements reproduce, without any
appeal to decision-theoretic intuition, the frozen-environment counterfactual that \cri{3} uses.

\begin{theorem}[Transport and grounded equivalence]\label{thm:transport}
$\{\ax{CF2},\ax{CF4},\ax{CF5}\}$ is minimally inconsistent.
\end{theorem}

\begin{proof}
Let $u(a,x)=x$, and compare $c(z)=z$ with $c'(z)=0$ at the baseline $0$. The instantiated component
outputs agree, so $E\sim_0E'$, but $V_E(1)=1$ and $V_{E'}(1)=0$. \ax{CF2} defines both values,
\ax{CF4} makes them unequal, and \ax{CF5} makes them equal. Omitting \ax{CF2}, \ax{CF4}, or
\ax{CF5}, the remaining pair is satisfied by $\kappa_{\mathrm{base}}$, $\kappa_{\mathrm{port}}$, or
$\kappa_{\mathrm{slot}}$ respectively.
\end{proof}

No computability or undecidability enters; the obstruction is the loss of information under an
equivalence relation. The second theorem connects that loss with the instability of
Part~\ref{part:single}.

\begin{theorem}[Calibration and ratifiability]\label{thm:calibration}
$\{\ax{CF3},\ax{CF5},\ax{CF10}\}$ is minimally inconsistent, and each of its proper subsets is satisfied
by a total, computable kernel satisfying \ax{CF1}, \ax{CF2}, \ax{CF7}, and \ax{CF8}.
\end{theorem}

\begin{proof}
Take the copied-action environment $u(a,x)=\one[a\neq x]$, $c(z)=z$. For either baseline,
$W_{E,b}(b)=0$ and $W_{E,b}(1-b)=1$. By Lemma~\ref{lem:projection} \ax{CF3} and \ax{CF5} force this table,
so every actual action recommends its opposite, contradicting \ax{CF10}. Omitting \ax{CF3}, \ax{CF5},
or \ax{CF10}, the remaining pair is satisfied by $\kappa_{\mathrm{flat}}$, $\kappa_{\mathrm{slot}}$, or
$\kappa_{\mathrm{port}}$, all total and computable.
\end{proof}

The witness environment is Death in Damascus with an infallible predictor
(Proposition~\ref{prop:did}). That no pure action is ratifiable there is familiar; what the theorem
adds is a derivation of the causal kernel from two semantic requirements that make no reference to
causation, and the exact minimality of the triple. Unlike the within-problem maps of
Section~\ref{sec:commit}, it is a minimal inconsistent family of three individually satisfiable
conditions.

\subsection{The complete map}

\begin{theorem}[Consistency map]\label{thm:cfmap}
On the finite context domain, the inclusion-minimal inconsistent subsets of
$\{\ax{CF1},\dots,\ax{CF10}\}$ are exactly
\begin{gather*}
  \{\ax{CF5},\ax{CF9}\},\quad
  \{\ax{CF2},\ax{CF4},\ax{CF5}\},\quad
  \{\ax{CF3},\ax{CF4},\ax{CF5}\},\\
  \{\ax{CF3},\ax{CF5},\ax{CF10}\},\quad
  \{\ax{CF4},\ax{CF5},\ax{CF10}\},\quad
  \{\ax{CF5},\ax{CF6},\ax{CF10}\},
\end{gather*}
and the maximal consistent subsets are exactly
\begin{center}
\small
\begin{tabular}{@{}ll@{}}
  \toprule
  witness & maximal consistent subset \\
  \midrule
  $\kappa_{\mathrm{slot}}$ & all except \ax{CF5} \\
  $\kappa_{\mathrm{port}}$ & $\{\ax{CF1},\ax{CF2},\ax{CF3},\ax{CF5},\ax{CF6},\ax{CF7},\ax{CF8}\}$ \\
  $\kappa_{\mathrm{base}}$ & $\{\ax{CF1},\ax{CF4},\ax{CF5},\ax{CF6},\ax{CF7},\ax{CF8}\}$ \\
  $\kappa_{\mathrm{flat}}$ & $\{\ax{CF1},\ax{CF2},\ax{CF5},\ax{CF7},\ax{CF8},\ax{CF10}\}$ \\
  \bottomrule
\end{tabular}
\end{center}
\end{theorem}

\begin{proof}
\emph{Inconsistency.} $\{\ax{CF5},\ax{CF9}\}$: \ax{CF9} gives \ax{CF2} and \ax{CF4}, and
Theorem~\ref{thm:transport} applies. $\{\ax{CF3},\ax{CF4},\ax{CF5}\}$: in the environment of
Theorem~\ref{thm:transport}, \ax{CF3} sets the value of action $1$ in the constant-context environment
$E'$ to $0$, \ax{CF5} transfers it to $E$, and \ax{CF4} requires $1$. $\{\ax{CF4},\ax{CF5},\ax{CF10}\}$:
take $u(a,x)=x$, $c(z)=z$, and let \ax{CF10} select $b$. The alternative $1-b$ is defined; \ax{CF5}
transfers definedness to $E^b$, whose component is the constant $b$, and \ax{CF4} requires the values
$1-b$ at $E$ and $b$ at $E^b$, which \ax{CF5} makes equal. $\{\ax{CF5},\ax{CF6},\ax{CF10}\}$: by
Lemma~\ref{lem:projection}, defined values are those of $\kappa_{\mathrm{port}}$, and \ax{CF10} defines
both values at the selected baseline in the environment of Theorem~\ref{thm:calibration}. The other two
sets are Theorems~\ref{thm:transport} and~\ref{thm:calibration}.

\emph{Minimality.} Deleting the members of $\{\ax{CF3},\ax{CF4},\ax{CF5}\}$ in turn leaves pairs
satisfied by $\kappa_{\mathrm{base}}$, $\kappa_{\mathrm{port}}$, and $\kappa_{\mathrm{slot}}$; for
$\{\ax{CF4},\ax{CF5},\ax{CF10}\}$ by $\kappa_{\mathrm{flat}}$, $\kappa_{\mathrm{slot}}$, and
$\kappa_{\mathrm{base}}$; for $\{\ax{CF5},\ax{CF6},\ax{CF10}\}$ by $\kappa_{\mathrm{slot}}$,
$\kappa_{\mathrm{flat}}$, and $\kappa_{\mathrm{port}}$. \ax{CF5} and \ax{CF9} are satisfied singly by
$\kappa_{\mathrm{port}}$ and $\kappa_{\mathrm{slot}}$. The remaining two cases are in
Theorems~\ref{thm:transport} and~\ref{thm:calibration}.

\emph{Exhaustiveness.} Let $\Sigma$ contain none of the six sets. If $\ax{CF5}\notin\Sigma$,
$\kappa_{\mathrm{slot}}$ witnesses $\Sigma$. Otherwise $\ax{CF9}\notin\Sigma$. If $\ax{CF3}\in\Sigma$, then
\ax{CF4} and \ax{CF10} are excluded and $\kappa_{\mathrm{port}}$ witnesses $\Sigma$. If $\ax{CF4}\in\Sigma$ and
$\ax{CF3}\notin\Sigma$, then \ax{CF2} and \ax{CF10} are excluded and $\kappa_{\mathrm{base}}$ witnesses
$\Sigma$. If neither \ax{CF3} nor \ax{CF4} is in $\Sigma$ but \ax{CF10} is, \ax{CF6} is excluded and
$\kappa_{\mathrm{flat}}$ witnesses $\Sigma$; if \ax{CF10} is also absent, $\kappa_{\mathrm{port}}$ does.
Each displayed set is maximal, since adding any further desideratum completes one of the six cores.
\end{proof}

Every core contains \ax{CF5}, and deleting \ax{CF5} leaves a single kernel satisfying everything else.

\begin{corollary}[A coherent semantics]\label{cor:parametric}
If \ax{CF5} is replaced by \emph{parametric extensionality}, $E\approx E'\Rightarrow\kappa(E,b,a)=\kappa(E',b,a)$,
then $\kappa_{\mathrm{slot}}$ satisfies all ten desiderata, and the agent choosing the first maximiser
of $V_E$ ratifies it.
\end{corollary}

\begin{proof}
$V_E$ depends on $E$ only through $u$ and the functions $c_i$, which parametric equivalence preserves;
the rest is Lemma~\ref{lem:witnesses}.
\end{proof}

The impossibility therefore concerns the erasure of dependency information, not extensionality in
every sense. The same positive model arises from a supplied dependency graph: a finite acyclic graph
of typed variables with total structural equations, a designated node for the decision computation,
and fixed exogenous data. Intervening on that node and recomputing its descendants, in the manner of
structural causal models \autocite{pearl2009}, retains the independent mechanisms and yields the
finite table $V_E$. The graphical formulation of functional decision theory intervenes in exactly
this way on a node representing the agent's computation \autocite[§§5--6]{yudkowskysoares2017}; what
the formulation does not provide, and what no axiom here supplies, is a procedure that extracts the
graph from an arbitrary closed program.

In matrix terms the map says the following. Grounded extensionality allows the kernel to see only
the host and the baseline outputs of the components, that is, the row of the causal matrix belonging
to the actual output, together with the payoff table. Calibration then pins the kernel to that row
(Lemma~\ref{lem:projection}), and Theorem~\ref{thm:calibration} is the observation that the diagonal
entry need not be maximal in its row. Parametric extensionality allows the kernel to see the whole
matrix, and then the diagonal is available. The choice between causal and ex-ante evaluation of
Part~\ref{part:single}, which there was a choice between two laws, reappears here as a choice between
two notions of when two computations are the same.

\subsection{Other conventions}\label{sec:otherconventions}

The map depends on how undefined values and realisability are treated, and these conventions are not
cosmetic. Suppose realisability is weakened to \ax{CF9$'$}: \emph{whenever} $\kappa(E,b,a)$ is defined,
it is attainable under slot substitution, so that \ax{CF9$'$} is equivalent to \ax{CF4}. Suppose
ratifiability is weakened to \ax{CF10$'$}: some computable $\tau$ has $\kappa(E,\tau(E),\tau(E))$ defined
and maximal among the defined values.

\begin{proposition}\label{prop:cfmapconditional}
With \ax{CF9$'$} and \ax{CF10$'$} in place of \ax{CF9} and \ax{CF10}, the minimal inconsistent subsets
are exactly
\begin{gather*}
  \{\ax{CF2},\ax{CF4},\ax{CF5}\},\quad \{\ax{CF2},\ax{CF5},\ax{CF9$'$}\},\quad
  \{\ax{CF3},\ax{CF4},\ax{CF5}\},\\ \{\ax{CF3},\ax{CF5},\ax{CF9$'$}\},\quad
  \{\ax{CF3},\ax{CF5},\ax{CF10$'$}\},\quad \{\ax{CF2},\ax{CF5},\ax{CF6},\ax{CF10$'$}\},
\end{gather*}
and the maximal consistent subsets are those of Theorem~\ref{thm:cfmap} with the primed axioms
substituted, except that $\kappa_{\mathrm{base}}$ now witnesses
\[
  \{\ax{CF1},\ax{CF4},\ax{CF5},\ax{CF6},\ax{CF7},\ax{CF8},\ax{CF9$'$},\ax{CF10$'$}\}.
\]
\end{proposition}

\begin{proof}
$\kappa_{\mathrm{base}}$ satisfies \ax{CF9$'$} and, with any computable $\tau$, \ax{CF10$'$}, so the cores
$\{\ax{CF4},\ax{CF5},\ax{CF10}\}$ and $\{\ax{CF5},\ax{CF6},\ax{CF10}\}$ do not survive. The cores containing
\ax{CF4} are duplicated with \ax{CF9$'$}, which is equivalent to it. \ax{CF3} and \ax{CF5} make $\kappa$
total and equal to $\kappa_{\mathrm{port}}$, which fails \ax{CF10$'$} in the copied-action environment.
\ax{CF2}, \ax{CF5}, and \ax{CF6} likewise give $\kappa_{\mathrm{port}}$ by Lemma~\ref{lem:projection}. The
four-element core is minimal: deleting \ax{CF2}, \ax{CF5}, \ax{CF6}, or \ax{CF10$'$} leaves a set satisfied
by $\kappa_{\mathrm{base}}$, $\kappa_{\mathrm{slot}}$, $\kappa_{\mathrm{flat}}$, or $\kappa_{\mathrm{port}}$.
For exhaustiveness, let $\Sigma$ be core-free and contain \ax{CF5} (otherwise $\kappa_{\mathrm{slot}}$
witnesses it). If $\ax{CF3}\in\Sigma$, $\kappa_{\mathrm{port}}$ witnesses it. If \ax{CF4} or \ax{CF9$'$} is in
$\Sigma$, \ax{CF2} is excluded and $\kappa_{\mathrm{base}}$ witnesses it. Otherwise, if \ax{CF10$'$} is in
$\Sigma$, either \ax{CF6} is absent and $\kappa_{\mathrm{flat}}$ witnesses $\Sigma$, or \ax{CF2} is absent
and $\kappa_{\mathrm{base}}$ does; if \ax{CF10$'$} is absent, $\kappa_{\mathrm{port}}$ does.
\end{proof}

The four-element core reflects the splitting of \enquote{defined and calibrated} into separate
requirements. Theorem~\ref{thm:calibration} is less sensitive to such bookkeeping, since its witnesses
are all total.

The word \enquote{realisable} also needs a benchmark. If realisability freezes every embedded
component at its reference implementation, an alternative taking $a$ obtains $W_{E,b}(a)$, and full
realisability in that sense forces $\kappa=\kappa_{\mathrm{port}}$; slot realisability forces
$\kappa=\kappa_{\mathrm{slot}}$. Neither follows from the word, and the benchmark must be supplied.

\subsection{Program domains and global totality}\label{sec:programmap}

The finite theorem extends to a program domain. Let the components be programs $c_i(d,a)\in B_i$,
promised to terminate on every string $d$ and every legal $a$, with a finite rational host, and with the
environment's description supplied as $d$. A fixed algorithm is allowed as a component $c(d,a)=F(d)$,
and it is independent in the declared sense, being invariant in the live parameter, even when $F$
happens to be the agent under evaluation. When \ax{CF10} asks only for a promise-total $\tau$, the four
kernels remain promise-computable, the slot optimiser halts on every valid description, and
Theorem~\ref{thm:cfmap} holds unchanged; no algorithm deciding membership in the promise class is
needed. Requiring $\tau$ to be globally total changes the map.

\begin{proposition}\label{prop:globalmap}
On the program domain, with \ax{CF10} requiring a globally total $\tau$, each of
$\{\ax{CF3},\ax{CF10}\}$, $\{\ax{CF4},\ax{CF10}\}$, $\{\ax{CF6},\ax{CF10}\}$, and $\{\ax{CF9},\ax{CF10}\}$ is
inconsistent. The minimal inconsistent subsets are these four together with $\{\ax{CF5},\ax{CF9}\}$,
$\{\ax{CF2},\ax{CF4},\ax{CF5}\}$, and $\{\ax{CF3},\ax{CF4},\ax{CF5}\}$, and the maximal consistent subsets are
$\{\ax{CF1},\ax{CF2},\ax{CF3},\ax{CF4},\ax{CF6},\ax{CF7},\ax{CF8},\ax{CF9}\}$ (witnessed by
$\kappa_{\mathrm{slot}}$, without a globally total optimiser) and the sets of $\kappa_{\mathrm{port}}$,
$\kappa_{\mathrm{base}}$, and $\kappa_{\mathrm{flat}}$ in Theorem~\ref{thm:cfmap}.
\end{proposition}

\begin{proof}
For a globally total candidate $\tau$, the component $c(d,a)=\tau(d)$ is admissible. With the host
$u(a,x)=\one[a\ne x]$ and $b=\tau(d(E_\tau))$, the component is constant in the live parameter, so
\ax{CF3}, \ax{CF4}, \ax{CF6}, and \ax{CF9} each force the values $0$ at $b$ and $1$ at $1-b$, contradicting
\ax{CF10}. The finite-domain cores remain, and those containing \ax{CF10} are now subsumed. If \ax{CF10} is
in a core-free set, \ax{CF3}, \ax{CF4}, \ax{CF6}, and \ax{CF9} are excluded, and the set lies in that of
$\kappa_{\mathrm{flat}}$, which every globally total agent ratifies; without \ax{CF10}, the case analysis of
Theorem~\ref{thm:cfmap} applies.
\end{proof}

This is the fixed-copy diagonal of Theorem~\ref{thm:fixedcopy}, factored through the desiderata, and
there is no conflict with promise-total optimisation: embedding the promise-total optimiser as $F$
need not yield a component that terminates on every string, so the diagonal need not belong to the
domain. If \ax{CF3} uses a syntactically narrower independence class that excludes embedded agent
algorithms, $\{\ax{CF3},\ax{CF10}\}$ need not be inconsistent; the proposition uses invariance in the
live parameter.

\subsection{The candidate semantics located}\label{sec:candidates}

The semantics of Section~\ref{sec:inference} can now be placed.

\emph{Transport through compositions.} For postcomposition alone, $h_0\circ T=h_1\circ T$ means that $h_0$
and $h_1$ agree on the image of $T$, which determines their values at every legal action exactly when
$T$ is surjective. So the relevant property is surjectivity, not injectivity, and in any case no
function from infinitely many descriptions to a finite action set is injective. Arbitrary contexts
restore the ambiguity for every agent: with $t=T(d)$, the contexts $C_0[P](x)=0$ and
$C_1[P](x)=\one[P(d)\ne t]$ are both the constant $0$ at $P=T$ and differ under a change of the output at
$d$. This is the environment of Theorem~\ref{thm:transport} in program form. With tagged contexts,
conditional transport, and grounded extensionality the core is $\{\ax{CF2},\ax{CF4},\ax{CF5}\}$; with
transport that asserts existence it is $\{\ax{CF4},\ax{CF5}\}$; with parametric extensionality the
contexts are not identified and $\kappa_{\mathrm{slot}}$ works.

\emph{Local matching} records each component's actual queried inputs, selects those components that
agree with $T$ there, and freezes the selection. On an appropriate termination promise it satisfies
\ax{CF1}, \ax{CF2}, agent extensionality, implementation extensionality at fixed inputs, and
\ax{CF8}. In the environment of Theorem~\ref{thm:localmatch} it violates \ax{CF3} and \ax{CF6} under
fixed-component independence, transport for transformed copies, \ax{CF9} under the ordinary
fixed-component benchmark, and \ax{CF10}. That does not make $\{\ax{CF1},\ax{CF2},\ax{CF5},\ax{CF8},\ax{CF10}\}$
inconsistent, which $\kappa_{\mathrm{flat}}$ satisfies.

\emph{Global matching} misses $1-T$, which is never equal to $T$, while transport requires its value
to become $1-a$. Its kernel can also be uncomputable with every component total: with $T\equiv0$ and
the components $g_e$ of Theorem~\ref{thm:haltred}, the environment returning $g_e(d(E_e))$ has actual
payoff $0$, and its value at action $1$ is $1$ if $e$ never halts and $0$ if it halts.

\emph{Fresh rematching} violates \ax{CF7} (Proposition~\ref{prop:rematch}: the final payoff is $1$ where
direct intervention gives $0$), but no desideratum requires it; retaining the original matching set, or
a supplied graph, avoids the failure. Literal syntactic substitution needs a one-pass convention for
the same reason, since the replacement code $T^{d\leftarrow a}$ mentions $T$ away from $d$ and repeated
rescanning can expand without end.

\emph{Material implication} yields the relation $\mathcal V$ of Section~\ref{sec:logical}, not a
kernel. In an independent five-and-ten the attainable payoff for ten is ten, so a selector assigning
$0$ violates \ax{CF3} and \ax{CF9}, but a selector from the relation need not: an additional rule could
choose the intervention value. A kernel assigning $(5,0)$ to an agent that takes five ratifies that
agent, so the failure is one of calibration and realisability, not of \ax{CF10}.

\emph{Fixed copies} are transported only if the equivalence convention erases their role. If
\ax{CF5} identifies a designated copy with a fixed program computing the same closed function, the
transport owed to the first passes to the second, against conservativity at a non-actual action;
if extensionality preserves the open interface, nothing passes. Under the fixed-component
realisation benchmark the global fixed-copy diagonal makes that realisability inconsistent with
\ax{CF10} alone.

\emph{Guarded logical semantics} for terminating programs is $\kappa_{\mathrm{base}}$, by
Proposition~\ref{prop:guardedarith}: it ratifies no agent under \ax{CF10} on a binary problem, and
every agent under \ax{CF10$'$}.

\emph{Logical induction with exploration} computes, at a stage $n$, values such as
\[
  \kappa_n(a)=\frac1n\sum_{i=0}^{n-1}\mathbb P_n\bigl(\bar U>i/n\mid A_n=a\bigr),
\]
using the capped conditional of \textcite[Definitions~4.7.1 and~4.8.2]{garrabrant2016li}, with value $1$
when the conditioning probability is zero. Nothing in the logical induction criterion makes these
values calibrated at a given stage: setting every price at one stage to zero preserves the criterion
by finite-perturbation closure, the capped conditional then assigns $1$ to every conditioned
proposition, and a constant-zero payoff receives $\kappa_n(0)=\kappa_n(1)=1$, violating \ax{CF1},
\ax{CF3}, and \ax{CF9}. The scope of exploration is a further input: with $u(a,b)=a+2b$ at the baseline
$a=b=0$, exploration that changes only the port yields the values $(0,1)$, while exploration that
changes the port and a designated copied bit together yields $(0,3)$. Conditioning computes the
consequences of a supplied exploration mechanism; it does not choose the mechanism.

\subsection{Robustness}

\begin{center}
\small
\renewcommand{\arraystretch}{1.2}
\begin{tabularx}{\textwidth}{@{}>{\raggedright\arraybackslash}p{10em}>{\raggedright\arraybackslash}X@{}}
  \toprule
  variation & effect \\
  \midrule
  globally total or promise-total agents & the finite obstruction survives both; arbitrary fixed total
     algorithms add the pairwise cores of Proposition~\ref{prop:globalmap} \\
  fixed algorithms admitted & fixed-target diagonals return when alternatives leave the target fixed;
     they are not needed for the finite obstruction \\
  fixed algorithms excluded & requires a concrete grammar: declared agent interfaces, no unrestricted
     evaluation or source reflection, intervention-total contexts; \enquote{nothing derived from $T$}
     is not a definition, since every $g$ is $C[T]$ for a context ignoring its argument \\
  informative agents & injectivity is unavailable; surjectivity removes only the $h\circ T$ ambiguity \\
  distribution-valued kernels & exact transport can demand the distinct point masses $\delta_0$ and
     $\delta_1$ for one grounded representation; uncertainty does not remove the contradiction \\
  mixed agents & some ratifiability failures acquire mixed fixed points, depending on prediction and
     sampling semantics (Section~\ref{sec:mixedexist}); the transport contradiction is untouched \\
  policy-level interventions & every one-input witness embeds; finite policy tables support the slot
     theorem; infinite policy spaces need separate assumptions \\
  bounded or certified languages & can support a globally total optimiser with a default on invalid
     descriptions; the unrestricted fixed-copy diagonal need not belong to them \\
  \bottomrule
\end{tabularx}
\end{center}
